\mychapter{1}{Conclusión}
\markboth{CONCLUSIÓN}{}

Recapitulación de lo obtenido y valoración personal del trabajo. La conclusión de fondo se puede
enunciar en una frase: sobre esta tarea, el modelo fundacional va por delante de la política clásica
por un margen pequeño y difícil de acreditar estadísticamente, mientras que la variable que sí
domina el resultado —con dos contrastes por debajo de $p = 0{,}001$— es la disposición de la
válvula. Es un resultado menos vistoso que «el modelo grande gana», y bastante más útil para quien
tenga que decidir dónde invertir el esfuerzo.


\subsection*{Aportaciones realizadas}

Enumeración de lo aportado por el trabajo:

\begin{itemize}
	\item Una tarea de manipulación industrial instrumentada y reproducible sobre el humanoide
	Unitree G1, con criterio de éxito objetivo, poses medidas y variabilidad controlada.
	\item Un conjunto de cien demostraciones teleoperadas, con dos disposiciones de la válvula y
	fondo fotorrealista aplicado por reproducción.
	\item Un banco de comparación que sirve a dos familias de políticas a través de la misma
	interfaz, lo que permite atribuir las diferencias al método.
	\item Una comparación emparejada con el contraste estadístico apropiado, y la caracterización
	del coste de cada método en datos y en latencia.
	\item El hallazgo, negativo pero accionable, de que la limitación dominante en esta tarea no es
	la cantidad de demostraciones sino su composición.
\end{itemize}

\subsection*{Trabajos futuros}

Líneas de continuación: reequilibrar el conjunto de datos hacia la disposición difícil, repetir los
entrenamientos con varias semillas para separar el efecto de los datos de la varianza del
entrenamiento, eliminar las dimensiones degeneradas del espacio de acción, ampliar a más de una
tarea y, sobre todo, validar sobre robot físico.

\subsection*{Problemas encontrados}

Relato de los obstáculos que condicionaron el desarrollo y de cómo se resolvieron. Son en buena
medida los que explican el reparto real del esfuerzo, muy escorado hacia la construcción de la
cadena frente al entrenamiento de los modelos.

\section*{Opiniones personales}

Valoración personal del trabajo: qué se ha aprendido, qué se haría de otra manera y qué percepción
deja el estado actual de los modelos fundacionales para robótica frente a las alternativas
específicas.
